\bookchapter{Глава 69. Справочник и ресурсы}

\chapterintro{Эта глава — справочник, которым можно пользоваться и отдельно. Забыли, как что-то пишется, или надо быстро принять решение — листайте прямо сюда. Держите его как «карманные карточки»: шпаргалка по ходовому коду, частые вопросы, ресурсы для учёбы — без единой пропущенной страницы.}

\booksection{Карточки быстрой справки по ходовому коду}

\begin{center}
\begin{tabularx}{\linewidth}{|X|X|}
\hline
\textbf{Что делаем} & \textbf{Ключевой код} \tabularnewline
\hline
Приватное состояние страницы & @State private var x = начальное значение \tabularnewline
Логика между страницами & \texttt{@Observable class VM \{ var x = 0 \}} \tabularnewline
Вертикальная / горизонтальная раскладка & \texttt{VStack \{ \}}/\texttt{HStack \{ \}} \tabularnewline
Гибкое пространство & \texttt{Spacer()} \tabularnewline
Список & List \{ ForEach(массив) \{ ... \} \} \tabularnewline
Переход на другую страницу & \texttt{NavigationStack + NavigationLink} \tabularnewline
Всплывающее окно снизу & .sheet(isPresented: \$show) \{ содержимое \} \tabularnewline
Закрыть текущую страницу & \texttt{@Environment(\textbackslash{}.\allowbreak dismiss) var dismiss} \tabularnewline
Загрузить данные из сети & \texttt{try await URLSession.\allowbreak shared.\allowbreak data(from: url)} \tabularnewline
Разобрать JSON & \texttt{struct X: Codable}+\texttt{JSONDecoder().\allowbreak decode} \tabularnewline
Сохранить мелкую настройку & UserDefaults.standard.set(значение, forKey: "ключ") \tabularnewline
Сохранить структурированные данные & \texttt{@Model}+\texttt{modelContainer}+\texttt{@Query} \tabularnewline
Поддержка тёмного режима & Берите системные цвета (.primary / Color(.systemGray6)), не вшивайте свои \tabularnewline
Запустить тесты & CmdU \tabularnewline
\hline
\end{tabularx}
\end{center}

\begin{center}\small\color{gray}Карточки быстрой справки по разработке iOS (август 2026)\end{center}

\booksection{Частые вопросы (FAQ)}

В 1: У меня нет Mac — могу ли я изучать разработку под iOS?

Язык изучить можно, полноценная разработка — нет. Swift Playgrounds на iPad/iPhone покрывает весь материал главы 61; но Xcode живёт только на macOS, так что в итоге Mac необходим. При ограниченном бюджете можно на переходный период арендовать облачный Mac или взять б/у Mac mini, а о вложениях решить, когда освоите язык.

В 2: Сколько стоит научиться разработке под iOS?

Этап учёбы: 0 (Xcode бесплатен + бесплатный Apple ID для отладки на устройстве). Этап публикации: аккаунт разработчика 99 долларов в год (около 700 юаней). Плюс вложения в сам Mac — вот и все расходы. А по сравнению с отдачей — даже маленькая утилита в App Store находит платящих пользователей — это выгодная сделка.

В 3: Что учить первым — SwiftUI или UIKit?

SwiftUI. Новые проекты, новые уроки и новые вакансии 2026 года строятся вокруг SwiftUI, и вся книга пользуется именно им. UIKit — фреймворк прошлого поколения, он по-прежнему живёт во множестве старых проектов: когда освоите SwiftUI, доучите UIKit по необходимости (у Apple есть учебники по взаимодействию), вберётся быстро. Только не перепутайте порядок.

В 4: Легко ли найти работу разработчика iOS?

Рынок 2026 года: начальные позиции «умею рисовать интерфейсы» перегреты конкуренцией, но разработчики, которые «знают Swift + SwiftUI и имеют опубликованное приложение», по-прежнему нарасхват. Медианные зарплаты iOS-разработчиков стабильно в верхних строчках среди стеков. Лучшая часть вашего резюме в период учёбы — то приложение в магазине, которое вы сделали сами.

В 5: Приложение отклонили на проверке — что делать?

Не паникуйте, это нормально — вероятность отказа при первой подаче немаленькая. Причины отказа Apple обычно пишет очень конкретно (какая страница, какой пункт правил); исправьте по каждому пункту, ответьте с пояснением в App Store Connect и подайте заново. Три частые причины отказа: непонятно назначение разрешений, в интерфейсе есть баги, метаданные (скриншоты/описание) не соответствуют действительности. Проверьте эти три пункта заранее — сэкономите круг ожидания.

\booksection{Частые горячие клавиши Xcode: шпаргалка}

\begin{center}
\begin{tabularx}{\linewidth}{|X|X|}
\hline
\textbf{Действие} & \textbf{Горячие клавиши} \tabularnewline
\hline
Запуск / тесты & CmdR / CmdU \tabularnewline
Открыть предпросмотр на холсте & OptionCmd \tabularnewline
Быстро открыть файл & CmdO \tabularnewline
Поиск по всему проекту & CmdF \tabularnewline
Форматировать отступы & CtrlI \tabularnewline
Закомментировать / раскомментировать & Cmd/ \tabularnewline
Очистить кеш сборки & CmdK \tabularnewline
Перейти к определению & Cmd + щелчок \tabularnewline
Переименовать (безопасно по всему проекту) & Правый клик Refactor Rename \tabularnewline
\hline
\end{tabularx}
\end{center}

\begin{center}\small\color{gray}Частые горячие клавиши Xcode: шпаргалка\end{center}

\begin{tcbinfo}{Совет: при «мистике» компиляции сперва чистим кеш}
В «полтергейсте» вида «вчера работало, сегодня нет» в шести-семи случаях из десяти виноват перепутанный кеш сборки. CmdK — очистить и пересобрать: это лечит большую половину мистики. Не помогло — перезапустите Xcode (вместе с Mac), не помогло и это — «удалите папку DerivedData»: таков заветный «тройной приём», передаваемый из уст в уста разработчиками iOS.

\end{tcbinfo}

\booksection{Рекомендуемые ресурсы для учёбы}

\booksection{Рекомендуемые ресурсы для учёбы}

\begin{center}
\begin{tabularx}{\linewidth}{|X|X|X|}
\hline
\textbf{Тип} & \textbf{Ресурс} & \textbf{Почему стоит} \tabularnewline
\hline
Официальный & developer.apple.com & Официальная документация Apple: там и уроки SwiftUI/Swift, авторитет номер один \tabularnewline
Официальный & SwiftUI Tutorials от Apple & Официальный курс «за руку», максимальная дружелюбность к новичкам \tabularnewline
Официальный & Официальный сайт Swift (swift.org) & Официальная документация языка и онлайн-площадка для упражнений \tabularnewline
Сообщество & Hacking with Swift & Бесплатная библиотека уроков Paul Hudson, качество общепризнанно в первой лиге \tabularnewline
Сообщество & Stack Overflow (тег swiftui) & Первая станция за помощью при ошибках \tabularnewline
Видео & Официальный канал Apple Developer (видео WWDC) & Авторитетные разборы новых технологий каждый год \tabularnewline
Дизайн & Руководство по интерфейсам Apple (HIG) & Официальная библия интерфейсов «с яблочным вкусом» \tabularnewline
К книге & Часть 7 «Практика программирования с ИИ» (главы 37–47) & Полное искусство разработки iOS с помощью ИИ \tabularnewline
\hline
\end{tabularx}
\end{center}

\begin{center}\small\color{gray}Список ресурсов для изучения разработки iOS (август 2026)\end{center}

\booksection{Заключение: от первого приложения к N-му}

Оглянитесь назад: от «не знаю, что такое SwiftUI» — до установленного Xcode, освоенного Swift, интерфейса, нарисованного на SwiftUI, подключённой к приложению базы данных и сети, и дальше до тестов, упаковки и публикации — самый трудный отрезок пути «от нуля к единице» вы уже прошли.

Мир технологий меняется быстро: фреймворки обновляются, инструменты сменяются, ИИ забирает всё больше «клавиатурной» работы. Но одно не устареет никогда: ваша способность размытую идею разложить на ясные шаги и шаг за шагом воплотить её. Инструменты — усилители, мышление — фундамент. Усилители можно менять, фундамент не должен шататься.

Вы больше не просто «умеющий пользоваться телефоном». Вы можете превратить идеи из головы в настоящие приложения на сотнях миллионов iPhone. С сегодняшнего дня творите с клавиатурой — в этом романтика iOS-разработчика.

\booksection{Список литературы}

\begin{enumerate}
  \item Apple Inc. (2026). SwiftUI Framework Documentation. developer.apple.com.
  \item Apple Inc. (2026). Swift Programming Language Guide. swift.org.
  \item Apple Inc. (2026). SwiftData Framework Documentation. developer.apple.com.
  \item Apple Inc. (2026). App Store Review Guidelines. developer.apple.com/app-store/review.
  \item Apple Inc. (2026). Human Interface Guidelines (HIG). developer.apple.com/design.
  \item Apple Inc. (2026). Xcode 27 Release Notes and App Distribution Guide.
\end{enumerate}
